\chapter{Testing and Evaluation}\label{chap:testing}

\section{Introduction}

Verification asks whether the system does what the requirements say. For an emergency platform it must also show that the
rules that protect people hold under conditions that are easy to get wrong: concurrent assignments, retried requests,
forged callbacks, withdrawn consent and failing dependencies. This chapter describes the verification strategy, the
automated tests that exist in the repository, and the end-to-end scenarios that check the requirements.

\section{Verification Strategy}

Verification combines three layers (\cref{fig:test-layers}), applied in proportion to risk.

\begin{enumerate}
  \item \textbf{Automated unit tests} cover the pure rule classes and parsers, where a mistake would be most costly and
        a test is cheapest: caller trust, identity-verification policy, webhook verification, facility rules, message
        construction and release parsing.
  \item \textbf{Static analysis and build checks} run in continuous integration. The mobile app is analysed on every pull
        request and release; the AI service runs its tests and image build on every push; and the backend validates its
        schema against its entities every time it starts.
  \item \textbf{Scenario verification} exercises the end-to-end workflows against a running stack, using the scenarios of
        \cref{sec:scenarios}, each traced to the requirements it covers.
\end{enumerate}

\begin{figure}[!htbp]\centering
\begin{tikzpicture}[x=1cm,y=1cm,font=\sffamily\small]
  \fill[najdablue!85] (-4.4,0) -- (4.4,0) -- (3.3,1.6) -- (-3.3,1.6) -- cycle;
  \fill[najdablue!55] (-3.3,1.7) -- (3.3,1.7) -- (2.2,3.3) -- (-2.2,3.3) -- cycle;
  \fill[najdanavy] (-2.1,3.4) -- (2.1,3.4) -- (0,5.1) -- cycle;
  \node[white,align=center] at (0,0.8) {\bfseries Scenario verification\\[-1pt]\footnotesize end-to-end workflows};
  \node[white,align=center] at (0,2.5) {\bfseries Static analysis and CI\\[-1pt]\footnotesize every push and pull request};
  \node[white,align=center] at (0,3.95) {\bfseries Unit tests\\[-1pt]\footnotesize\textcolor{white}{pure rules}};
  \node[anchor=west,najdagrey,align=left,font=\sffamily\footnotesize] at (4.7,0.8) {few, broad, slow};
  \node[anchor=west,najdagrey,align=left,font=\sffamily\footnotesize] at (4.7,3.95) {many, narrow, fast};
\end{tikzpicture}
\caption{The verification layers, from fast unit tests to end-to-end scenarios.}\label{fig:test-layers}
\end{figure}

\section{Automated Tests in the Repository}

\Cref{tab:tests-auto} lists the automated tests present in the repository and what they protect. The backend suite
concentrates on the classes whose behaviour is decisive and deterministic; the mobile tests cover the offline outbox, the
Egypt boundary check and the data models; the AI service tests cover the model registry and the manifest contract.

{\tabfont\rowband
\begin{longtable}{@{}>{\raggedright\arraybackslash}p{\dimexpr(\linewidth-6\tabcolsep)*197/1000\relax}>{\raggedright\arraybackslash}p{\dimexpr(\linewidth-6\tabcolsep)*421/1000\relax}>{\raggedright\arraybackslash}p{\dimexpr(\linewidth-6\tabcolsep)*382/1000\relax}@{}}
\caption{Automated tests in the repository}\label{tab:tests-auto}\\
\rowcolor{najdanavy}\textcolor{white}{\sffamily\bfseries\small Area} & \textcolor{white}{\sffamily\bfseries\small Test classes /\allowbreak{} files (number of tests)} & \textcolor{white}{\sffamily\bfseries\small What they protect}\\
\endfirsthead
\rowcolor{najdanavy}\textcolor{white}{\sffamily\bfseries\small Area} & \textcolor{white}{\sffamily\bfseries\small Test classes /\allowbreak{} files (number of tests)} & \textcolor{white}{\sffamily\bfseries\small What they protect}\\
\endhead
\bottomrule\endlastfoot
\textbf{Backend --- caller trust} & CallerTrustPolicyTest (8) & Weights, decay, window, tier thresholds and cool-down effect.\\
\textbf{Backend --- identity verification} & KycPolicyTest (5), KycStatusTest (3), KycVendorDataTest (3), DiditWebhookVerifierTest (6), DiditClientListTest (3), DiditClientDeleteTest (3) & Who must verify, status mapping from provider labels, vendor-data format, signature and timestamp verification, provider client behaviour.\\
\textbf{Backend --- facility management} & FacilityManagerPolicyTest (5), FacilityCsvParserTest (5), FacilityCsvSeederTest (5) & Facility-type rules for units and staff; parsing and idempotent seeding of the bundled facility list.\\
\textbf{Backend --- incident evidence} & IncidentEvidenceTest (6) & Handling of incident evidence in the AI triage path.\\
\textbf{Backend --- notifications} & NotificationMessagesTest (5) & Bilingual message construction for every notification type.\\
\textbf{Backend --- releases} & AppReleaseServiceTest (6), ReleaseVersionsTest (4) & Release lookup, caching and version comparison.\\
\textbf{Backend --- application} & NajdaApplicationTests (1) & The application context starts and the Flyway-managed schema matches the entities.\\
\textbf{Mobile (Flutter), 40 tests} & outbox\_\allowbreak{}report\_\allowbreak{}test, geo\_\allowbreak{}boundary\_\allowbreak{}test, medical\_\allowbreak{}profile\_\allowbreak{}test, notifications\_\allowbreak{}test, push\_\allowbreak{}preference\_\allowbreak{}test and further core and feature tests & Offline outbox behaviour, Egypt boundary check, medical-profile model, notification model, push preference.\\
\textbf{AI service (pytest), 2 tests} & test\_\allowbreak{}registry & Model registry, manifest contract and export validation.\\
\textbf{Total} & 68 backend + 40 mobile + 2 AI service & 110 automated tests.\\
\end{longtable}}

\section{Verification Scenarios}\label{sec:scenarios}

The scenarios in \cref{tab:tests-scenarios} check the requirements end to end. Each states an action, the expected
result and the requirements it covers. They are written to be executed against a running stack (the Compose environment
of \cref{chap:implementation}) by a tester who plays the roles involved. The scenarios concentrate on the behaviours that are
easy to get wrong: atomic registration, idempotent submission, race conditions on unit assignment, consent and
privacy rules, abuse control, after-commit notifications and failure of external services.

{\tabfont\rowband
\begin{longtable}{@{}>{\raggedright\arraybackslash}p{\dimexpr(\linewidth-8\tabcolsep)*89/1000\relax}>{\raggedright\arraybackslash}p{\dimexpr(\linewidth-8\tabcolsep)*380/1000\relax}>{\raggedright\arraybackslash}p{\dimexpr(\linewidth-8\tabcolsep)*367/1000\relax}>{\raggedright\arraybackslash}p{\dimexpr(\linewidth-8\tabcolsep)*165/1000\relax}@{}}
\caption{Verification scenarios}\label{tab:tests-scenarios}\\
\rowcolor{najdanavy}\textcolor{white}{\sffamily\bfseries\small ID} & \textcolor{white}{\sffamily\bfseries\small Scenario} & \textcolor{white}{\sffamily\bfseries\small Expected result} & \textcolor{white}{\sffamily\bfseries\small Covers}\\
\endfirsthead
\rowcolor{najdanavy}\textcolor{white}{\sffamily\bfseries\small ID} & \textcolor{white}{\sffamily\bfseries\small Scenario} & \textcolor{white}{\sffamily\bfseries\small Expected result} & \textcolor{white}{\sffamily\bfseries\small Covers}\\
\endhead
\bottomrule\endlastfoot
\textbf{TC-01} & Register with e-mail: request code, enter correct code, complete the form. & Account created, signed in, profile complete, terms recorded. & FR-IAM-01--04, UC-1\\
\textbf{TC-02} & Enter a wrong code five times, then request another within 60 s. & Attempts exhausted message; resend refused until cool-down ends. & FR-IAM-01, BR-2\\
\textbf{TC-03} & Force a database failure after the Firebase account is created. & The Firebase account is deleted; a retry succeeds. & FR-IAM-03\\
\textbf{TC-04} & Provision a dispatcher as Admin; open the one-time link; sign in. & Dispatcher workspace opens; no password was shown. & FR-IAM-12, UC-12\\
\textbf{TC-05} & Plain Admin tries to change a Super Admin’s role, or to disable the last Super Admin. & Refused with an explanation. & FR-IAM-14, FR-IAM-16\\
\textbf{TC-06} & Complete Didit verification; replay the same webhook; send a webhook with a bad signature. & Status VERIFIED once; replay ignored; bad signature rejected. & FR-ID-01--03\\
\textbf{TC-07} & Publish a Terms version that requires acceptance; call a protected endpoint. & 403 LEGAL\_\allowbreak{}CONSENT\_\allowbreak{}REQUIRED until accepted; users notified. & FR-LEG-01--02, UC-57\\
\textbf{TC-08} & Submit a report with GPS; repeat the same request id. & One incident; dispatchers notified once; the repeat returns the same incident. & FR-REP-01, FR-REP-03\\
\textbf{TC-09} & Submit with coordinates outside Egypt, and a manual pin 1 km from the device. & Both refused with the reason. & FR-REP-02, BR-16\\
\textbf{TC-10} & Submit a second report while one is active. & 409 conflict. & FR-REP-04\\
\textbf{TC-11} & Submit offline; restore connectivity. & Report arrives automatically exactly once. & FR-REP-13, UC-28\\
\textbf{TC-12} & Report for someone else with sharing on; read the incident’s medical profile as a dispatcher. & Profile not shared. & FR-MED-04\\
\textbf{TC-13} & Withdraw sharing consent while an incident is open; open the profile as the crew. & Not shared immediately; no audit entry for a read. & FR-MED-03, FR-MED-05\\
\textbf{TC-14} & Open a patient profile as an Admin, then as a crew member not on the mission. & Both refused. & FR-MED-05, UC-47\\
\textbf{TC-15} & Two dispatchers assign the same unit at the same moment. & One mission is created; the other gets a conflict. & FR-DSP-06\\
\textbf{TC-16} & Assign a unit with no LEAD on shift. & Refused. & FR-DSP-06\\
\textbf{TC-17} & LEAD accepts, goes en route, arrives, completes (fire truck). & Citizen sees each step; unit AVAILABLE; incident RESOLVED. & FR-FLD-05--07\\
\textbf{TC-18} & Cancel the last active mission while the incident is ASSIGNED. & Incident returns to DISPATCHER\_\allowbreak{}REVIEW. & FR-DSP-08\\
\textbf{TC-19} & Citizen cancels after a unit arrived. & 409; response continues. & FR-REP-08\\
\textbf{TC-20} & Mark a false report before and after a unit has arrived. & Closed in the first case; verdict-only in the second; audit entry both times. & FR-DSP-10\\
\textbf{TC-21} & Accumulate false reports for one caller; submit again. & 429 with Retry-After and the emergency numbers; reporting works after the wait. & FR-DSP-12\\
\textbf{TC-22} & Ambulance: select hospital, en route, send vitals, arrive. & Hospital alerted; vitals streamed live; mission completed; incident resolved. & FR-HOS-01--05\\
\textbf{TC-23} & Send vitals for a free-text destination. & Refused. & FR-HOS-04, FR-HOS-06\\
\textbf{TC-24} & Facility manager creates a unit and a staff member at a fire station. & Fire truck and firefighter only; other roles refused. & FR-ADM-08\\
\textbf{TC-25} & Delete a unit with history; one on a live mission. & Archived; refused. & FR-ADM-05\\
\textbf{TC-26} & Roll back a transaction that would have notified (for example a failing assignment). & No notification is created. & FR-NOT-03\\
\textbf{TC-27} & Send an urgent notification to a device with an unregistered token. & Push fails once, token deleted, inbox row remains. & FR-NOT-05\\
\textbf{TC-28} & Switch to Arabic on web and mobile. & Right-to-left layout and Arabic notification text. & FR-PLT-01, FR-NOT-02\\
\textbf{TC-29} & Stop the language-model provider and submit a report. & Report created; no priority; dispatcher can retry later. & FR-DSP-02\\
\textbf{TC-30} & Open the status endpoint with the database stopped. & 503 with database DOWN. & FR-PLT-02\\
\end{longtable}}


\section{Test Environment and Procedure}

Scenarios are run against the local stack: PostgreSQL, the backend, the AI service and the web console started by Docker
Compose, with a Firebase project and test accounts for each role, and the mobile app installed on an Android device or
emulator. Failure scenarios are produced by stopping or misconfiguring the relevant dependency (for example, removing the
language-model key or stopping the database). Each scenario is recorded with its identifier, the observed result and a
pass or fail verdict. The automated suites are run with the standard tool of each component; the backend suite includes a
test that starts the whole application, so it needs a reachable PostgreSQL database.

\section{Results}\label{sec:results}

\Cref{tab:results} gives the outcome of each suite together with the command that reproduces it.

\begin{table}[!htbp]
\caption{Verification results.}\label{tab:results}
\centering\tabfont\rowband
\begin{tabular}{@{}>{\bfseries\raggedright\arraybackslash}p{4.0cm}>{\raggedright\arraybackslash}p{4.2cm}>{\raggedright\arraybackslash}p{5.8cm}@{}}
\rowcolor{najdanavy}\textcolor{white}{\sffamily\bfseries\small Suite} & \textcolor{white}{\sffamily\bfseries\small Command} & \textcolor{white}{\sffamily\bfseries\small Result} \\
Backend tests (15 classes) & \code{./mvnw test} in \code{backend/} & 68 run; 0 failures, 0 errors, 0 skipped \\
Mobile tests & \code{flutter test} in \code{mobile/} & 40 run; all passed \\
AI service tests & \code{pytest -q} in \code{ai/} & 2 run; all passed \\
Verification scenarios & TC-01 to TC-30 (\cref{tab:tests-scenarios}) & 30 of 30 passed \\
\end{tabular}
\end{table}

All 110 automated tests and all 30 verification scenarios passed. The backend run includes the application-start test,
which applies the Flyway migration to PostgreSQL and lets Hibernate validate the schema against every entity, so it
also confirms that the database and the code agree. The results describe the code base as delivered; because each suite
is a single command, and the scenarios are written down, they can be repeated after any change.

\section{Discussion}

The verification strategy deliberately puts the most effort where the risk is highest. The pure rule classes are
tested exhaustively because they are cheap to test and costly to get wrong; the concurrency, idempotency and
after-commit guarantees are checked by dedicated scenarios because they cannot be demonstrated by a unit test of a single
class. Two limitations should be stated openly: the repository's automated coverage is concentrated on rule classes
rather than on whole services, and the end-to-end scenarios are executed by hand rather than by a browser-automation
suite. Both are addressed in the future work (\cref{sec:future-work}).

\section{Chapter Summary}

This chapter set out a risk-proportionate verification strategy, inventoried the automated tests, and defined 30
end-to-end scenarios traced to the requirements. The next chapter assesses what has been achieved against the objectives.